\documentclass[11pt]{amsart}
\usepackage[T1]{fontenc}
\usepackage{lmodern}
\usepackage{amssymb}
\usepackage[hidelinks]{hyperref}

\hypersetup{pdftitle={Polynomial rigidity for anisotropic minimal graphs in dimensions four and five},pdfauthor={Kaiwen Jin}}
\newtheorem{theorem}{Theorem}
\allowdisplaybreaks[2]
\urlstyle{same}

\raggedbottom
\begin{document}

\title[Polynomial rigidity for anisotropic minimal graphs]{Polynomial rigidity for anisotropic minimal graphs in dimensions four and five}
\author{Kaiwen Jin}
\email{yui.ui.siki@gmail.com}
\keywords{Anisotropic minimal graphs, Bernstein problem, polynomial solutions, Jacobi fields}

\begin{abstract}
An entire anisotropic minimal graph over \(\mathbb R^4\) or \(\mathbb R^5\) defined by a polynomial must be affine if the surface integrand is positive, even, \(C^{2,1}\) on the sphere, and uniformly elliptic in tangential directions. The boundaries of the superlevel sets of the leading homogeneous term minimize a restricted surface energy. A comparison of positive Jacobi fields reduces the five-dimensional case to an indefinite quadratic factor. When its cofactor is nonconstant, the Jacobi equation fixes a noninteger scaling exponent for a cluster of polynomial roots. Analyticity of the cluster's centroid then gives the contradiction.
\end{abstract}

\maketitle

\section{Introduction}

The anisotropic Bernstein problem asks when an entire stationary graph for a uniformly elliptic surface energy must be a hyperplane. Mooney and Yang constructed nonpolynomial examples over \(\mathbb R^4\) with subquadratic growth \cite{MY}. Over \(\mathbb R^6\), Mooney constructed a quadratic example for a \(C^{2,1}\) surface integrand \cite[Theorem 1.1]{Mooney6}. The polynomial question in the intervening dimensions is posed in \cite[Section 7.2]{MooneySurvey}. We prove rigidity in both dimensions for the full uniformly elliptic \(C^{2,1}\) class of positive even integrands.

\begin{theorem}\label{thm:main}
Let \(n\in\{4,5\}\) and let \(\Phi\) be a positive even \(C^{2,1}\) function on the unit sphere \(S^n\subset\mathbb R^{n+1}\). Use the same symbol for its positively one-homogeneous extension. Suppose that there is a \(\lambda>0\) such that
\[
 D^2\Phi(\nu)[v,v]\ge \lambda |v|^2
 \quad\text{for all }\nu\in S^n,\ v\in\nu^\perp.
\]
If a real polynomial \(u\) on \(\mathbb R^n\) satisfies
\[
 \operatorname{div} DF(Du)=0\quad\text{on }\mathbb R^n,\qquad
 F(p)=\Phi((-p,1)),
\]
then \(u\) is affine.
\end{theorem}

Chen and Lan prove rigidity for entire smooth solutions through
base dimension seven under a \(C^2\) proximity condition to the
Euclidean area integrand \cite[Theorem~1.1]{CL}. For smooth uniformly
elliptic integrands, Yang proves rigidity if the gradient grows as
\(o(R^\alpha)\) on balls of radius \(R\), for an exponent \(\alpha>0\)
depending on dimension and integrand \cite[Theorem~1.1]{Yang}.

The proof uses translated dilations to extract the levels of the
leading homogeneous term and then considers their normal velocities,
as in Guo's treatment of Euclidean area
\cite[Theorem~4.7 in the preprint]{Guo}. We use the regularity of
elliptic minimizing boundaries \cite[Part~I]{SSA} and compare
positive homogeneous Jacobi fields for the anisotropic operator.
In dimension four, the boundary of the positive phase of the leading
term must be a hyperplane. In dimension five, the leading term has
an indefinite quadratic factor. A constant cofactor gives a
hyperbolic energy equation, as in \cite[Remark~3.5]{Mooney6}. For a
nonconstant cofactor we retain the lower-order terms and compare the
polynomial on transverse lines to a Jacobi field on the product of
the quadratic cone and the vertical line.

Suppose that \(D=\deg u\ge2\), and let \(U\) denote the homogeneous
term of degree \(D\). Write \(G(p)=\Phi((p,0))\) for the integrand
on the horizontal hyperplane, where the last coordinate is zero.

\section{Calibrated graphs, minimizing levels, and regular convergence}\label{sec:blowdown}

The energy of an oriented hypersurface for an integrand \(\Psi\)
is the surface integral of \(\Psi(\nu)\), with \(\nu\) the unit
normal. For a set \(E\), we use \(\partial E\) for its oriented
boundary current in comparisons and use the support of that current
as its boundary surface.

The homogeneity identity \(D^2\Phi(\xi)\xi=0\) and the positive tangential
Hessian make the extension convex, with value zero at the origin. In
particular, the supporting-plane inequality
\(D\Phi(\xi)\cdot v\le\Phi(v)\) holds for \(\xi\ne0\). The vector field
\[
 X(x,z)=D\Phi((-Du(x),1))
\]
has zero ambient divergence by the graph equation. Its pairing with the
upward unit normal to the graph equals the graph energy density, by
Euler's identity. The associated closed form certifies minimization
against compactly supported integral current variations. Indeed the
pairing integrates to the same value under a compactly supported
boundary variation by Stokes, and is bounded above by the energy of
each competitor. The upper bound applies to either orientation.
Dilations and translations preserve this minimizing property.

We use Schoen and Simon's results for locally minimizing oriented
boundaries of sets in ambient dimension at least three \cite[Part~I]{SSA}.
At boundary points, their volume estimate (35) gives uniform two-phase
density, and (33) gives local mass bounds. With these mass bounds and
local-volume convergence of the sets, Theorem 1.1(vi)--(vii) preserves
minimization and gives energy convergence in almost every interior comparison
ball. Theorem 1.2 and Remarks 1 and 3 give local single-graph
convergence through second derivatives at interior regular points of the limit.

Each integrand used below is autonomous and \(C^{2,1}\) on the sphere. By
multiplying by a positive constant we may arrange the lower bounds in
(7)--(8) of \cite[Part~I]{SSA}. The \(C^{2,1}\) bound on the
sphere and homogeneity give the bounds in (9), including the bound
for weak third derivatives in the normal variable. All
spatial derivatives vanish. The same verification applies to
restrictions to linear subspaces. In the local comparison class
of \cite[Part~I, Section~1.1]{SSA}, the difference \(S'-S\) between a
same-boundary integer-rectifiable competitor \(S'\) and the restricted
boundary current \(S\) is an integral cycle in the comparison ball.
Coning gives a filling \(B\) supported in the ball with
\(S'-S=\partial B\). We use the two-dimensional regularity of
\cite[Part~I, Corollary~3.2 and the following remark]{SSA}, in its
form for absolutely minimizing integral currents. We restrict the
minimizing boundary to an interior ball at a radius at which the restriction
is integral and its boundary is supported on the sphere. Competitors with
the same boundary supported in the closed ball are admissible compact
variations of the full surface. Such radii have full measure by slicing,
and the radius condition on the integrand is automatic in the autonomous
case.

For each \(c\in\mathbb R\), the subgraphs of
\(R^{-1}u(Rx)-R^{D-1}c\) converge in local volume to
\(\{U>c\}\times\mathbb R\). Indeed the term of order \(R^{D-1}\) is
\(U-c\), and the exceptional level has zero Lebesgue measure. The limit
therefore minimizes \(\Phi\). Its base boundary \(\partial\{U>c\}\)
minimizes \(G\). To see the latter claim, if a bounded base perturbation
reduces the \(G\)-energy by \(\delta>0\), insert that perturbation in the
cylinder for vertical coordinates in \((-T,T)\) and restore the
original cylinder above and below it. The saving on the sides is
\(2T\delta\), and the two caps cost a bounded amount independent of
\(T\). For a current variation \(\partial B\) of the base boundary,
the variation in the cylinder is \(\partial(B\times[-T,T])\).
The cap energies are bounded by \(2\sup_{S^n}\Phi\,\mathbf M(B)\),
where \(\mathbf M(B)\) denotes the mass of the filling. The side energy
is computed with \(G\) because the normals are horizontal. This proves
minimization in the required current class.

We record two facts for subsequent use. First, \(U\) changes sign.
Otherwise the unshifted graph blow-down subgraphs tend to a constant
set, in contradiction with two-phase density at \((0,R^{-1}u(0))\).
Second, for any nonzero positive-degree homogeneous polynomial \(P\)
whose superlevel boundaries all minimize an autonomous elliptic energy,
\begin{equation}\label{eq:zeros}
 \{P=0\}=\operatorname{spt}\partial\{P>0\}.
\end{equation}
The inclusion from right to left is immediate. If \(P(x)=0\) and \(x\)
were outside the support, one phase of \(\{P>0\}\) would have zero volume
in a ball about \(x\). Choose \(x_k\to x\) with \(P(x_k)=c_k\ne0\).
Euler's formula makes \(x_k\) a smooth boundary point of
\(\{P>c_k\}\). The two-phase density bound there contradicts the
convergence of the latter sets to a constant phase in the ball.

Finally, if the zero boundary of \(U\) were a hyperplane \(H\), the
unshifted dilated graphs would converge regularly near the origin to
the vertical hyperplane \(H\times\mathbb R\). Their normals at
\((0,R^{-1}u(0))\) are independent of \(R\) and have a strictly positive
upward component. A normal to the limiting hyperplane has upward
component zero. Regular convergence gives a contradiction.

\section{Positive homogeneous Jacobi fields on a regular cone}\label{sec:exponents}

For an autonomous integrand \(\Psi\) and a stationary hypersurface \(M\),
write \(A=D_M\nu\) for its shape operator, \(P_\Psi=D^2\Psi(\nu)|_{TM}\),
and \(V=\operatorname{tr}(P_\Psi A^2)\). The weak Jacobi equation is
\begin{equation}\label{eq:jacobi}
 \mathcal J_\Psi f=\operatorname{div}_M(P_\Psi\nabla_M f)+V f=0.
\end{equation}
This operator governs the normal velocities of stationary families,
including translations. For the regularity used here it is enough to
derive the form by normal variations. In the variation \(x\mapsto x+t f\nu\),
the integrand times the area element, relative to the original area
element, is
\(\det B\,\Psi(\nu-t B^{-1}\nabla_M f)\), where \(B=I+t f A\).
Writing \(H=\operatorname{tr} A\) and \(Y=(D\Psi(\nu))_T\), the second
derivative at zero is
\[
 \langle P_\Psi\nabla_M f,\nabla_M f\rangle+
 D_M(f^2)\cdot (A-HI)Y+f^2\Psi(\nu)(H^2-\operatorname{tr} A^2).
\]
Codazzi gives \(\operatorname{div}_M(A-HI)=0\), and
\(D_M Y=P_\Psi A-\Psi(\nu)A\). Taking the trace of
\((A-HI)D_M Y\) gives
\[
 \operatorname{div}_M((A-HI)Y)
 =V-H\operatorname{tr}(P_\Psi A)+\Psi(\nu)(H^2-\operatorname{tr}A^2).
\]
The stationarity condition is \(\operatorname{tr}(P_\Psi A)=0\).
Integration by parts now gives the quadratic form
\(\int_M(\langle P_\Psi\nabla_M f,\nabla_M f\rangle-V f^2)\).
Polarization and differentiation of stationary first variation give
\eqref{eq:jacobi} for the normal velocity. The expansion uses only derivatives of
\(\Psi\) through order two. All hypersurfaces where we apply \eqref{eq:jacobi}
are smooth locally as polynomial graphs, regular algebraic level
sets, or their cylindrical products.

Consider a connected component of a regular stationary cone away
from the origin in \(\mathbb R^d\). Its intersection \(S\) with the
unit sphere is the link. Here \(S\) is compact and has no boundary.
At \(\theta\in S\), split \(P_\Psi\) relative to the radial direction
and \(T_\theta S\) as
\(\left(\begin{smallmatrix}a&b^{\!T}\\b&B\end{smallmatrix}\right)\).
Write \(V(r\theta)=r^{-2}v(\theta)\), and \(s=d-3\). In the cone
metric the angular equation for \(r^\ell\psi(\theta)\) is
\begin{equation}\label{eq:angular}
 \mathcal A_\ell\psi=
 \operatorname{div}_S(B\nabla_S\psi)+(s+2\ell)b\cdot\nabla_S\psi+
  [\ell\operatorname{div}_S b+\ell(s+\ell)a+v]\psi=0.
\end{equation}
Indeed the radial divergence contributes the factor \(\ell+s\) to
the radial component of the flux, and the tangential divergence
contributes \(\ell\operatorname{div}_S(b\psi)\). Integration by parts
gives \(\mathcal A_\ell^*=\mathcal A_{-s-\ell}\).

A smooth positive solution \(\psi\) of \(\mathcal A_\ell\psi=0\)
yields a positive weak solution at the paired exponent \(-s-\ell\).
For detail, \(Lh=\psi^{-1}\mathcal A_\ell(\psi h)\) is a uniformly
elliptic divergence form operator with Lipschitz principal coefficients
and drift, and no zeroth-order term. It has a positive invariant
density \(\rho\). To construct it, large \(c-L\) is invertible
\(H^1\to H^{-1}\) by coercivity. The embedding through \(L^2\) is compact,
so \(L\) and its adjoint are Fredholm of index zero. The compact
maximum principle \cite[Theorems~8.19--8.20]{GT} gives \(\ker L=\mathbb R\). Any nonzero element of
the adjoint kernel has nonzero integral, since otherwise \(Lh=1\)
would be solvable, contrary to the same principle. Normalize its
integral to one. For smooth \(g\ge0\), Fredholm compatibility solves
\(Lh=g-\int_S \rho g\). If the integral were negative, the right-hand
side would be strictly positive everywhere, which is impossible on
the compact link. Thus \(\rho\ge0\) and local Harnack applied to its
adjoint equation gives \(\rho>0\). Lipschitz coefficients and drift
give the usual weak elliptic regularity and bounded zeroth-order
coefficient for this adjoint equation. It follows by testing
\(\mathcal A_\ell(\psi h)\) that \(\mathcal A_\ell^*(\rho/\psi)=0\).

There cannot be three distinct real exponents with positive link
solutions. If \(\ell_1<\ell_2<\ell_3\), choose \(t\in(0,1)\) with
\(\ell_2=t\ell_1+(1-t)\ell_3\). Applying \eqref{eq:jacobi} to the geometric mean
of the corresponding outer positive fields gives, for
\(z=\psi_1^t\psi_3^{1-t}\) and
\(\xi=(\ell_1-\ell_3,\nabla_S\log\psi_1-\nabla_S\log\psi_3)\),
\begin{equation}\label{eq:mean}
 \frac{\mathcal A_{\ell_2}z}{z}
  =-t(1-t)\langle P_\Psi\xi,\xi\rangle<0.
\end{equation}
Weak link solutions are locally \(C^{1,\beta}\) for some \(\beta>0\).
Indeed, after interior \(W^{2,2}\) regularity, expansion in charts gives
a strong equation with continuous principal coefficients and bounded
first- and zeroth-order coefficients. Since the link has dimension
at most three, Sobolev embedding gives local \(L^4\) bounds for the
solution and its first derivatives. Moving the lower-order terms
to the right then gives an \(L^4\) right-hand side. Interior
higher integrability gives \(W^{2,4}_{\mathrm{loc}}\), which embeds
into \(C^{1,\beta}_{\mathrm{loc}}\) \cite[Theorems~8.8 and 9.11 and
Lemma~9.16]{GT}. The weak
product computation therefore uses admissible products on compact
charts, with each positive solution bounded away from zero. The inequality is strict uniformly on
the link by ellipticity and \(\ell_1\ne\ell_3\). Dividing by the strictly
positive solution at \(\ell_2\) gives a strict supersolution for
an elliptic operator without a zeroth-order term on a compact
boundaryless manifold, contradicting the minimum principle.

If the cone is nonflat, any smooth positive homogeneous Jacobi
field satisfies \(-s<\ell<0\). For \(\ell<-s\) or \(\ell>0\), mix
\(\ell\) and its paired exponent to obtain \eqref{eq:mean} at zero. Since
\(v\ge0\), the resulting strict supersolution also contradicts the
minimum principle for \(\mathcal A_0-v\). If \(\ell=0\) or \(-s\),
there is a positive solution at zero. It is a supersolution for
\(\mathcal A_0-v\), hence constant by the strong minimum principle.
It follows that \(v=0\) identically, and ellipticity gives \(A=0\),
contrary to nonflatness.

\section{Structure of the leading term in dimensions four and five}\label{sec:leading}

\subsection{A factor on a regular component}

Let \(P\) be a nonzero homogeneous polynomial of positive degree
whose superlevel boundaries all minimize an autonomous positive
elliptic \(C^{2,1}\) energy. Assume its zero boundary is regular away from
the vertex. At a generic point of any connected component of that
regular cone, precisely one real irreducible homogeneous factor \(p\)
vanishes and \(Dp\ne0\). One way to see genericity is to choose a
transverse direction where \(P\) does not vanish, and project the patch
along that direction. The square-free part of \(P\), with its leading
constant scaled out, is monic in the transverse variable. It is
square-free also over the rational function field of the other
variables by Gauss's lemma. The nonzero discriminant in those variables
then gives simple transverse roots on a dense subset of the open
projection. Write \(P=p^m W\) there, with \(W\) the polynomial
quotient. Since \(P\) changes sign across its minimizing boundary
and \(W\ne0\) locally, \(m\) is odd.

In fact this factor and multiplicity persist all along the component,
with \(Dp\ne0\) and \(W\ne0\). The points with these properties are
relatively open. At any limit point in the regular component the
factor \(p\) still vanishes. In a fixed graph chart near that point,
the levels \(\{P=c\}\) for small positive \(c\) converge as single
graphs to \(\{P=0\}\), by \eqref{eq:zeros}, local volume convergence and the
minimizing-boundary regularity quoted above. Their graph difference
has a strict sign by nesting and disjointness. In fixed graph
coordinates the positive difference solves a uniformly elliptic
divergence-form equation with coefficients obtained by averaging the
graph integrand Hessian along the two slopes. Harnack compares its
values at the limit point and at a point of the open factor patch
inside the chart. The displacement on the open patch is comparable
to \(c^{1/m}\) since the transverse derivative of \(p\) is nonzero
and \(W\ne0\). The displacement magnitude \(h_c\) at the limit point
is therefore also comparable to \(c^{1/m}\). If \(Dp=0\) there,
then \(p=o(h_c)\) on the displaced graph. If \(W=0\) there, then
\(W=o(1)\) and \(p=O(h_c)\) on the displaced graph. In either
case \(c=p^m W=o(h_c^m)=o(c)\), which proves relative closedness.

The local signed root \(p W^{1/m}\) is smooth along this component
and homogeneous of degree \(\deg(P)/m\). Nearby levels of the root
are stationary for the same energy, including the zero level.
The positive normal velocity is \(1/|D(pW^{1/m})|\). Its homogeneity
is \(\ell=1-\deg(P)/m\). Every connected link component is compact,
since it is a connected component of the closed zero link, which
is regular. If a component is nonflat, Section~\ref{sec:exponents} gives
\begin{equation}\label{eq:degree-bounds}
 1<\deg(P)/m<d-2,
\end{equation}
where \(d\) is the ambient base dimension. A linear factor at a
regular component gives a hyperplane component, because \(Dp,W\)
remain nonzero there and the component is both open and closed
in the corresponding punctured hyperplane.

\subsection{Reduction at a nonzero boundary point}

The regularity of the cone away from its vertex follows by dimension reduction. For a nonzero
point \(x_0\) of the zero boundary of \(P\), let \(T\) be the first
nonzero polynomial in the Taylor expansion of \(P\) at \(x_0\), and
let \(k\ge1\) be its degree. Comparing the terms of order \(k-1\) in
Euler's identity gives \(x_0\cdot DT=0\). The sets
\(\{y:P(x_0+\varepsilon y)>\varepsilon^k c\}\) converge in local
volume to \(\{T>c\}\) for all \(c\), so the corresponding superlevel
boundaries are minimizing. By two-phase density at the old zero boundary point,
the limit at \(c=0\) has both phases. The limit is a cylinder in
the \(x_0\)-direction, and the side-and-cap argument of Section~\ref{sec:blowdown}
gives minimizing superlevel boundaries of \(T|_{x_0^\perp}\) in dimension \(d-1\)
for the restricted integrand. If \(d=4\), the reduced zero boundary
is a two-dimensional minimizing cone. It is regular at its vertex
by \cite[Part~I, Corollary~3.2 and the following remark]{SSA}. A cone regular at its vertex is a
hyperplane, as seen by comparing its tangent plane at zero with
its unchanged rescalings. The limiting cylinder in dimension four
is therefore planar, and regular convergence makes \(x_0\) regular
for the original zero boundary.

In dimension four, \eqref{eq:degree-bounds} would force \(\deg(P)/m<2\), whereas a
nonflat component has \(\deg p\ge2\) and \(\deg(P)/m\ge\deg p\).
Thus all components are planar. Each such component is a whole
punctured hyperplane, by relative openness, closedness and
connectedness. Distinct hyperplanes intersect away from zero and
cannot be distinct regular components. We have proved that in
dimension four a homogeneous polynomial with minimizing superlevel
boundaries and a sign-changing zero boundary has just one
hyperplane zero boundary. Applied to \(P=U\), this contradicts
the fixed graph normal in Section~\ref{sec:blowdown}.

In dimension five, at any nonzero zero boundary point the same
Taylor argument gives a four-dimensional homogeneous polynomial
with minimizing superlevel boundaries and a sign-changing zero boundary.
The preceding four-dimensional statement makes the reduced
boundary planar. Regular convergence again gives regularity of
the original five-dimensional zero boundary away from the vertex.

\subsection{The quadratic in dimension five}

For \(P=U\) in dimension five, there must be a nonflat component
by Section~\ref{sec:blowdown}. By \eqref{eq:degree-bounds} its irreducible factor has degree two. Such
a quadratic has both signs since it has regular real zeros. If its
positive and negative inertias are \(a,b\), the Hessian of the
quadratic restricted to the tangent space at a regular zero has
\(a-1\) positive and \(b-1\) negative eigenvalues. For example, use
the null point and a companion null vector with nonzero quadratic
pairing to split off one hyperbolic plane. The remainder of the
tangent quadratic form is the orthogonal complement of that plane,
with the original kernel and the radial direction contributing
zeros. If \(\min(a,b)=1\), this second fundamental form is semidefinite
and nonzero, unless \(a=b=1\), which would factor over the reals.
The contraction with the positive tangential energy Hessian at a
regular point cannot vanish for a nonzero semidefinite form.
Thus \(a,b\ge2\).

After an invertible linear coordinate change, the regular zero set
of the quadratic is parametrized by the
product of the two positive-dimensional unit spheres, a common
positive radius, and a free kernel direction. It is therefore connected.
The zero-boundary
component carrying this factor is open and closed within that regular
zero set, by its local single-factor description and its closedness
in the punctured zero boundary. If the quadratic had a nontrivial
kernel, the component would accumulate at a nonzero kernel point.
Closedness would put that point in the component, contradicting
the established nonvanishing of \(Dp\) there. Thus its rank is
five, of inertia \((3,2)\) or \((2,3)\).

Two distinct nonflat components would have distinct quadratic
factors by the connectedness just proved. Their multiplicities
\(m_1,m_2\) would each exceed \(D/3\) by \eqref{eq:degree-bounds}, whereas their total
degree \(2(m_1+m_2)\) is at most \(D\). A hyperplane component
would intersect the quadratic cone. Indeed, on any hyperplane
the quadratic takes both signs, since each of its sign subspaces
has dimension at least two. It therefore has a nonzero zero in
that hyperplane. The regular components are disjoint. This also
shows that the cofactor has no other real zeros away from the
origin, by \eqref{eq:zeros} and its nonvanishing along the quadratic component.

We conclude
\begin{equation}\label{eq:factorization}
 \begin{gathered}
 U=q^m W,\qquad m\text{ odd},\qquad j=\deg W\text{ even},\\
 0\le j<m,\qquad D=2m+j,\qquad W>0\text{ off }0.
 \end{gathered}
\end{equation}
Here \(q\) is full-rank indefinite of inertia \((3,2)\) up to sign.
Change \(q\) to \(-q\) if necessary to arrange the sign of \(W\).
Homogeneity and positivity imply the asserted parity of \(j\).
The bound \(j<m\) is exactly \(D/m<3\).

\section{A constant cofactor}\label{sec:constant}

If \(j=0\) then all levels of \(H=qW^{1/m}\) are stationary for \(G\).
Its gradient is an invertible linear map. The trace equation for
its levels gives \(D^2G(p):D^2H=0\) for \(p\ne0\), first away from
the zero level and then by continuity. Set \(M=D^2H\), and write
\(p=L\widetilde p\) and \(\widetilde G(\widetilde p)=G(L\widetilde p)\)
for an invertible matrix \(L\). The transformed equation has
coefficient \(L^{-1}ML^{-T}\) against \(D^2\widetilde G\).
By the inertia of \(M\), choose \(L\) so that the coefficient is
\(I_3\oplus(-I_2)\) up to a common sign. The transformed energy is
one-homogeneous and \(C^{2,1}\) away from the origin. Average it
over \(O(3)\times O(2)\), preserving the linear equation and its
positivity. With the two blocks denoted \(p_+,p_-\), write its average
as \(|p_+|h(v)\) for \(v=|p_-|/|p_+|\). Then
\begin{equation}\label{eq:ode}
 (v^2-1)h''-(2v+1/v)h'+2h=0.
\end{equation}
This follows by applying \(\Delta_{p_+}-\Delta_{p_-}\), both in
radial coordinates. Put \(h_1(v)=1+v^2/2\), a solution of \eqref{eq:ode}.
On either side of \(v=1\) reduction of order gives
\[
 (h/h_1)'=c_\pm\,\frac{|v^2-1|^{3/2}}{v h_1(v)^2}.
\]
At \(v=0\) a nonzero \(c_-\) gives logarithmic divergence, so
\(h=h(0)h_1\) for \(v<1\). At \(v=1^+\), a nonzero \(c_+\) gives
a correction with nonzero leading term proportional to
\((v-1)^{5/2}\). Its second derivative has leading order
\((v-1)^{1/2}\), which violates \(C^{2,1}\). Thus continuity at
one forces \(h=h(0)h_1\) for all \(v\). As \(p_+\to0\) with
\(p_-\ne0\) fixed, the resulting energy diverges, contrary to
continuity away from the origin. This excludes the constant-cofactor case.

\section{Rescaled vertical derivatives on the cylinder}\label{sec:derivatives}

Assume \(j\ge2\) in \eqref{eq:factorization}. Put \(C=\{q=0\}\setminus\{0\}\) and
\(\Gamma=C\times\mathbb R\), the regular part of the unshifted
blow-down surface. Use \(\nu_C=Dq/|Dq|\) on \(C\), pointing toward
\(q>0\). For \(y\in C\) and \(z\in\mathbb R\), define for large \(R\)
\begin{equation}\label{eq:graph}
 P_{R,y}(t)=R^{-1}u(R(y+t\nu_C(y))),\qquad
 P_{R,y}(t_R(y,z))=z,
\end{equation}
where \(t_R\) is the small normal displacement of the rescaled graph
over \((y,z)\in\Gamma\). Section~\ref{sec:blowdown} gives a single normal graph with
\(C^2\) convergence to zero displacement locally on \(\Gamma\).
Its local charts agree on overlaps by graph uniqueness. In
particular the slope \(P_{R,y}'(t_R)\) is nonzero. To see its positivity, observe that the oriented normal \((Du,-1)/\sqrt{1+|Du|^2}\) of the graph tends
to \((\nu_C,0)\) by convergence to the boundary of \(\{q>0\}\times
\mathbb R\) with orientation towards \(\{q>0\}\). This orientation
can be used since \(\Phi\) is even. Let
\[
 J_R(y,z)=\partial_z t_R(y,z)=(P_{R,y}'(t_R(y,z)))^{-1}>0.
\]
If \(\nu_R\) is this oriented normal, then
\begin{equation}\label{eq:velocity}
 -\nu_R\cdot e_{n+1}
    =(\nu_R\cdot(\nu_C,0))J_R.
\end{equation}
This follows by differentiating the normal graph in \(z\) and
taking its tangent derivative's normal component. The factor on
the right tends to one uniformly on compact patches.

The left side of \eqref{eq:velocity} is the normal velocity of vertical translation
(up to sign), hence solves \eqref{eq:jacobi} on the rescaled graph. Under the
normal graph maps to \(\Gamma\), the Jacobi equation has uniformly
elliptic divergence coefficients and bounded zeroth-order terms
on compact subcharts. The coefficients converge to those of \eqref{eq:jacobi}
on \(\Gamma\) by \(C^2\) convergence. Interior Harnack applied
along a finite chain of such subcharts compares values of \(J_R\)
on each compact subset of the connected cylinder with its value
at one fixed reference point \((y_*,0)\). The constants for a
given compact set are uniform in \(R\). The same local equations
give uniform H\"older estimates and local \(H^1\) bounds for the
normal components after division by the reference value,
by interior elliptic estimates and the Caccioppoli inequality.
Using \eqref{eq:velocity}, local uniform limits of the normalized \(J_R\) exist
along subsequences, stay strictly positive, and satisfy the
cylindrical Jacobi equation in its weak form. The weak passage
uses the \(H^1\) bounds for the normal components and the
convergence of coefficients.

\subsection{The normalization scale}

The scale is determined by an algebraic branch.
If \(a(s)>0\) is a continuous algebraic branch for \(s>0\) small,
then
\begin{equation}\label{eq:algebraic}
 a(s)\sim c s^\alpha,\qquad c>0,\quad \alpha\in\mathbb Q.
\end{equation}
For completeness, write a nonzero polynomial relation as
\(\sum_i p_i(s)a(s)^i=0\), omitting identically zero \(p_i\).
Every remaining \(p_i\) has an integer vanishing order \(v_i\).
At each sufficiently small positive \(s\), the term largest in
absolute value must be comparable within a fixed factor to
another term, by the equality. This bounds \(a(s)\) between
\(c_1s^\gamma\) and \(c_2s^\gamma\) for one of the finitely many
rational slopes \(\gamma=(v_i-v_k)/(k-i)\). For small \(s\)
the bands for different slopes are disjoint. Continuity on a
connected interval selects one band, so \(a(s)/s^\gamma\) stays
in a fixed compact subset of \((0,\infty)\). Divide the relation
by the smallest power among \(s^{v_i+i\gamma}\). Every limit
of \(a(s)/s^\gamma\) is a root of the resulting nonzero limiting
polynomial. Its cluster set is connected, being the intersection
of the nested closed images of shrinking intervals, and is
contained in a finite root set. Thus it consists of a single
positive \(c\), proving \eqref{eq:algebraic}.

For each fixed \((y,z)\), set \(s=1/R\) in \eqref{eq:graph} to get
\[
 s^D u((y+t\nu_C(y))/s)-s^{D-1}z=0
 \qquad\text{at }t=t_R(y,z).
\]
This is a polynomial relation in \(s,t\) because \(u\) has degree \(D\).
The reciprocal derivative formula consequently makes \(J_R\) algebraic
in \(s\).
The small root, and hence \(J_R\), depends continuously on \(s\)
by transversality and local uniqueness for all sufficiently
small positive \(s\). At the reference point choose \(\alpha\)
from \eqref{eq:algebraic}, so \(J_R(y_*,0)\sim c R^{-\alpha}\). Harnack
comparison then forces the same \(\alpha\) at every fixed
\((y,z)\), by \eqref{eq:algebraic} applied there. Combining pointwise uniqueness
with the compactness above gives
\begin{equation}\label{eq:limit}
 R^\alpha J_R\longrightarrow J>0
 \quad\text{locally uniformly on }\Gamma,\qquad
 \mathcal J_\Phi J=0\text{ on }\Gamma.
\end{equation}
This step only concerns compact subsets of the regular cylinder.

\subsection{The limiting inverse polynomial}

Let \(T_R(y)=t_R(y,0)\) and introduce the centered variables
\begin{equation}\label{eq:centered}
 w_R(y,z)=R^\alpha(t_R(y,z)-T_R(y)),\qquad
 f_{R,y}(w)=P_{R,y}(T_R(y)+R^{-\alpha}w).
\end{equation}
By \eqref{eq:limit}, \(w_R\) tends locally uniformly to
\(w(y,z)=\int_0^z J(y,h)\,dh\). For each \(y\), evaluate the
polynomial \(f_{R,y}\) at \(D+1\) distinct chosen height nodes
\(w_R(y,z_i)\). Their values are exactly \(z_i\). The nodes
converge to distinct nodes because \(J>0\). Inverting the
Vandermonde matrix shows coefficientwise convergence to a
polynomial \(f_y\) of degree at most \(D\), whose coefficients
depend continuously on \(y\). The convergence is uniform on compact \(y\)-patches, on which the
limiting node separations have a positive minimum. The identities pass to the limit
for every real \(z\), giving \(f_y(w(y,z))=z\).
As \(z\to\pm\infty\), \(w(y,z)\) has infinite limits with
the same signs, since it is increasing and the polynomial
is bounded on any compact interval of \(w\). Thus \(f_y\)
is a bijection of \(\mathbb R\) with
\(f_y'(w(y,z))J(y,z)=1\). In particular \(f_y'>0\) on
the entire real line.

The coefficient of \(w^m\) in \(f_{R,y}\) is
\[
 R^{D-1-\alpha m}\left[
 \frac{1}{m!}\partial_t^m U(y+T_R(y)\nu_C(y))+O(R^{-1})
 \right],
\]
where the \(t\)-derivative denotes differentiation in the displayed
normal direction. The bracket has a nonzero limit because
\(q\) has nonzero normal derivative at \(y\) and \(W(y)>0\).
Boundedness of this coefficient gives
\begin{equation}\label{eq:alpha-bound}
 \alpha\ge(D-1)/m=2+(j-1)/m>2.
\end{equation}
The same coefficient formula shows that every power higher
than \(m\) vanishes in the limit, since its prefactor has a
strictly negative exponent and the bracket stays bounded.
Consequently \(1\le\deg f_y\le m\) for every \(y\).

\section{Cone exponents and the degree of the limiting polynomial}\label{sec:degree}

The cylindrical Jacobi coefficients are independent of \(z\).
A Harnack chain on the unit link of \(C\) therefore compares
\(J(y,z)\) at any two unit link points with the same height,
by constants uniform over all \(z\). The link is connected
by the quadratic description in Section~\ref{sec:leading}. If \(K_y=\deg f_y\)
and \(a_{K_y}(y)\) is its positive leading coefficient,
inverting the polynomial in Section~\ref{sec:derivatives} gives
\begin{equation}\label{eq:inverse}
 J(y,z)\sim
 \frac{1}{K_y}a_{K_y}(y)^{-1/K_y}z^{-1+1/K_y}
 \quad(z\to+\infty).
\end{equation}
Indeed \(w(y,z)\sim (z/a_{K_y}(y))^{1/K_y}\) and
\(J=1/f_y'(w)\). Distinct \(K_y\) at two points would
make their ratio unbounded in one direction. Hence there
is a single degree \(K\) on the unit link. For \(r>0\),
the identity
\[
 t_R(ry,z)=r\,t_{Rr}(y,z/r)
\]
follows from \eqref{eq:graph}, since \(\nu_C(ry)=\nu_C(y)\).
Differentiating in \(z\) and using \eqref{eq:limit} gives
\(J(ry,z)=r^{-\alpha}J(y,z/r)\). Thus also
\(w(ry,z)=r^{1-\alpha}w(y,z/r)\) and
\(f_{ry}(w)=r f_y(r^{\alpha-1}w)\). The degree is \(K\)
on all of \(C\), and
\begin{equation}\label{eq:radial-law}
 a_K(ry)=r^{1+(\alpha-1)K}a_K(y).
\end{equation}

On every compact subset of \(C\), the leading coefficient
is bounded below by a positive constant, and all coefficients
are bounded. Therefore the elementary inversion in \eqref{eq:inverse} is
uniform on such sets. For \(Z\to+\infty\) the functions
\(Z^{1-1/K}J(y,Z+z)\) converge locally uniformly on
\(\Gamma\) to
\[
 A_*(y)=K^{-1}a_K(y)^{-1/K}>0 .
\]
Each is a cylindrical Jacobi field by vertical translation
invariance. Caccioppoli and weak passage on each compact regular
patch show that the limit solves \eqref{eq:jacobi} there. It is independent
of \(z\). Test the weak equation with \(\varphi(y)\eta(z)\), where
both factors have compact support and \(\int_{\mathbb R}\eta=1\).
The mixed horizontal--vertical coefficient contributes a multiple
of \(\int_{\mathbb R}\eta'=0\), since the coefficients and \(A_*\)
are independent of \(z\). The horizontal tangential block is
\(D^2G(\nu_C)|_{TC}\), and the zero vertical principal curvature
reduces the potential to its base expression. The remaining weak
equation is the \(G\)-Jacobi equation on \(C\). By \eqref{eq:radial-law} its
degree is \(\ell_*=-\alpha+1-1/K<-1\).

On the same cone the smooth positive normal speed of the
stationary levels of \(qW^{1/m}\), from Section~\ref{sec:leading}, has degree
\(\ell_-=1-D/m<-1\). In base dimension five the paired degree for a positive
field of Section~\ref{sec:exponents} is \(\ell_+=-2-\ell_->-1\).
If \(\ell_*\ne\ell_-\), there would be three distinct real
degrees with positive solutions on the compact regular link.
Section~\ref{sec:exponents} rules this out. Therefore
\[
 -\alpha+1-1/K=1-D/m,\qquad
 \alpha=D/m-1/K.
\]
Combining this with \eqref{eq:alpha-bound} and \(K\le m\) forces \(K=m\) and
\(\alpha=(D-1)/m\).

\section{The centroid of the root cluster}\label{sec:cluster}

Fix one \(y\in C\). With \(s=R^{-1}\) put
\[
 Q(s,t)=s^D u((y+t\nu_C(y))/s)
       =\sum_{d=0}^D s^{D-d} U_d(y+t\nu_C(y)),
\]
where \(U_d\) are the homogeneous parts of \(u\).
This is a polynomial in \(s,t\). At \(s=0\), zero has
multiplicity exactly \(m\) as a root in \(t\). Choose a
small circle about zero enclosing no other root of \(Q(0,t)\).
For all small complex \(s\) there are \(m\) roots in the
circle, counted with multiplicity. Their centroid
\[
 c(s)=\frac{1}{2\pi i m}\oint t\,\frac{Q_t(s,t)}{Q(s,t)}\,dt
\]
is analytic and is real for real \(s\) by complex conjugation.
Here the circle is positively oriented.

For positive \(s\) we have exact equality
\[
 f_{R,y}(w)=s^{-(D-1)}Q(s,T_R(y)+s^\alpha w).
\]
The left side converges coefficientwise to \(f_y\), now
known to have degree exactly \(m\). On a fixed circle in
the \(w\) plane enclosing all its roots, Rouch\'e's theorem
puts \(m\) roots of \(f_{R,y}\) inside for all large \(R\).
Since \(T_R(y)\to0\), these correspond to precisely the
\(m\) roots in the fixed \(t\) circle. In particular
\(|c(s)-T_R(y)|=O(s^\alpha)\). Write \(b_s=(c(s)-T_R(y))/s^\alpha\)
and choose a subsequence \(s\downarrow0\) for which \(b_s\to b\in\mathbb R\).
The change of center is
\[
 s^{-(D-1)}Q(s,c(s)+s^\alpha w)=f_{R,y}(w+b_s).
\]
Along the subsequence the polynomial on the right converges coefficientwise
to \(f_y(w+b)\), whose linear coefficient is \(f_y'(b)>0\).
The prelimit linear coefficient is
\[
 s^{\alpha-(D-1)}\,Q_t(s,c(s)).
\]
The analytic factor either has an integer vanishing order or vanishes
identically. Thus a finite nonzero subsequential
limit requires \(D-1-\alpha\) to be an integer. This is
impossible because \(\alpha=2+(j-1)/m\) and \(2\le j<m\).

The contradiction excludes every nonconstant \(W\) in \eqref{eq:factorization}
with all lower homogeneous terms of \(u\) allowed. Section~\ref{sec:constant}
excludes the constant cofactor. Together with Section~\ref{sec:leading},
this proves the stated affine rigidity for both dimensions.

\subsection*{Acknowledgment}
ChatGPT/Codex assisted with the development and checking of the argument and with the preparation of the manuscript.

\begin{thebibliography}{99}

\bibitem{CL}
L. Chen and J. Lan,
\emph{A \(C^2\)-perturbative Bernstein theorem for anisotropic entire minimal graphs},
preprint, 2026.
\href{https://arxiv.org/abs/2608.29117v2}{arXiv:2608.29117v2}.

\bibitem{GT}
D. Gilbarg and N. S. Trudinger,
\emph{Elliptic Partial Differential Equations of Second Order},
Classics in Mathematics, Springer, Berlin, 2001.
\href{https://doi.org/10.1007/978-3-642-61798-0}{doi:10.1007/978-3-642-61798-0}.

\bibitem{Guo}
Y. Guo,
\emph{On polynomial solutions to the minimal surface equation},
Calc. Var. Partial Differential Equations \textbf{65} (2026), article 141.
\href{https://doi.org/10.1007/s00526-026-03316-2}{doi:10.1007/s00526-026-03316-2}.
\href{https://arxiv.org/abs/2404.00115v2}{arXiv:2404.00115v2}.

\bibitem{Mooney6}
C. Mooney,
\emph{Entire solutions to equations of minimal surface type in six dimensions},
J. Eur. Math. Soc. \textbf{24} (2022), 4353--4361.
\href{https://doi.org/10.4171/JEMS/1202}{doi:10.4171/JEMS/1202}.

\bibitem{MooneySurvey}
C. Mooney,
\emph{Bernstein theorems for nonlinear geometric PDEs},
survey, 2024.
\href{https://arxiv.org/abs/2407.11903v2}{arXiv:2407.11903v2}.

\bibitem{MY}
C. Mooney and Y. Yang,
\emph{The anisotropic Bernstein problem},
Invent. Math. \textbf{235} (2024), 211--232.
\href{https://doi.org/10.1007/s00222-023-01222-4}{doi:10.1007/s00222-023-01222-4}.

\bibitem{SSA}
R. Schoen, L. Simon, and F. J. Almgren Jr.,
\emph{Regularity and singularity estimates on hypersurfaces minimizing parametric elliptic variational integrals},
Acta Math. \textbf{139} (1977), 217--265.
\href{https://doi.org/10.1007/BF02392238}{doi:10.1007/BF02392238}.

\bibitem{Yang}
Y. Yang,
\emph{A growth gap for anisotropic minimal graphs},
preprint, 2026.
\href{https://arxiv.org/abs/2609.08972v1}{arXiv:2609.08972v1}.

\end{thebibliography}

\end{document}
